\documentclass[pdflatex,sn-mathphys-num]{sn-jnl}
\usepackage{amsmath,amssymb,booktabs,array,microtype}
\usepackage{xurl}
\providecommand{\burl}[1]{\url{#1}}
\raggedbottom
\begin{document}
\title[Evidence selection in budgeted multi-hop RAG]{Evidence Selection in Budgeted Multi-hop Retrieval-Augmented Generation: From Proxy Objectives to Answer Quality}
\abstract{Evidence-set selection for multi-hop question answering must identify complementary passages within a limited reader context. We study whether supervised set scores remain useful on the sets selected by search. First-, second-, and fourth-order factorization models and a general permutation-invariant model share dense candidates, support-annotation supervision, and a fixed reader. Despite high accuracy on constructed set pairs, the static fourth-order model loses annotated support relative to a size-matched dense set that is itself available to the final selector. Search-aligned continuation then trains the deployed score on merged search frontiers and final sets, with an equal-update static replay control. On the same 128-question development panel, the primary fourth-order model's annotated support completeness rises from 55/128 to 57/128, while answer F1 falls from 0.4181 to 0.3766; maximal marginal relevance scores 0.4303. The second seed also shows higher support completeness and lower answer F1. Lower-order responses differ, including small gains over the corresponding static second-order model in both seeds. Exact granular-ball and k-means indexes preserve flat beam-search choices but do not accelerate the measured 128-candidate workload. These results locate a practical mismatch among constructed-set prediction, deployed selection, and reader utility, with a separate limit on indexing efficiency. They characterize exposed development conditions rather than independent confirmation; earlier sentence-completion results remain a distinct method and evaluation boundary.}
\keywords{Retrieval-augmented generation, Evidence selection, Multi-hop question answering, Set prediction, Search-aligned learning, Empirical evaluation}
\maketitle

\section{Introduction}
A retrieval-augmented question-answering system must decide which evidence
to place in its reader's context. For multi-hop questions, individual
relevance is only part of that decision: several passages may jointly supply
a connection that no passage provides alone. A large candidate list can
contain the required facts and still yield an incomplete or distracting
context after selection. Conversely, adding more retrieved material can
consume a limited input budget without improving an answer.

This motivates evidence-set models, which score combinations rather than
isolated passages. Recent retrieval work explicitly models collective
information needs~\cite{lee2025setr}, and multi-hop beam retrieval retains
multiple partial hypotheses~\cite{zhang2024beam}. Such designs raise a
measurement question: does a model that recognizes good evidence sets also
select them when optimized by its actual search procedure? The two tasks
need not induce the same input distribution. Constructed positive and
negative sets may be easy to distinguish, while search discovers sets that
exploit inaccurate scores. The resulting selection must then pass a second
test: helping a particular reader produce a correct answer.

We investigate these transitions in a controlled development study. Frozen
BGE embeddings retrieve a shared list of 128 paragraphs. Four compact
models receive the same known-support supervision: first-order (H1),
second-order (H2), fourth-order (H4), and a general DeepSets control.
All select under the same actual 1,024-token input limit, use the same
six-block maximum, and feed the same unadapted FP16 reader. The analysis
distinguishes sampled-set prediction, actual selected support, canonical
answer quality, and measured cost. A bounded continuation then changes the
training distribution and preference loss while retaining the model
families, candidates, renderer, and search.

The investigation originates in HyperGranular-RAG, a sentence-level
structured evidence-completion system. Its earlier compact-backbone gains
are real but conditional: they did not establish an advantage over a
stronger BGE retriever. The present paragraph-set models are a subsequent
design, not a corrected implementation of that historical heuristic.
Likewise, using granular balls to index conditional set scores changes
search computation, not the selected evidence when the index is exact.
Keeping these method identities separate prevents a later mechanism from
being used to explain an earlier score retrospectively.

The article contributes three controlled observations. A high-accuracy
set scorer can choose less annotated support even when a size-matched dense
alternative is included in its final decision. A search-aligned intervention
can improve annotated support completeness yet lower answer F1 on the same
questions, with distinct responses from low-order and general-set controls.
Finally, exact indexes can preserve those beam choices while increasing
latency at 128 candidates. These observations link the training target to
actual selection and reader behavior, and separate selection quality from
the cost of computing it.

\section{Related work and the scope of comparison}
\subsection{Retrieval, compression, and set selection}
Dense passage retrieval represents queries and passages in a shared space
for efficient relevance ranking~\cite{karpukhin2020dpr}. Multi-hop methods
such as MDR condition later retrieval on previously found evidence
\cite{xiong2021mdr}, while IRCoT alternates retrieval with generated reasoning
\cite{trivedi2023ircot}. Beam Retrieval jointly trains an encoder and
hop-dependent heads to preserve multiple passage hypotheses
\cite{zhang2024beam}. Our candidate universe is fixed before set selection;
we do not reproduce those systems or borrow their reported performance.

Compression and joint selection are also established research directions.
RECOMP trains extractive and abstractive compressors for downstream language
model utility~\cite{xu2024recomp}. SetR uses information-requirement reasoning
and teacher-generated set selections~\cite{lee2025setr}. Our models operate
on cached embeddings and known support annotations, without an answer-based
teacher or a generated summary. These differences define a low-cost
experimental boundary. They do not establish superiority to compression or
LLM-based set selection, which are not directly evaluated here.

MMR balances query relevance and redundancy~\cite{carbonell1998mmr}.
Coverage-based summarization objectives similarly formalize complementarity
\cite{lin2011submodular}. Both are important alternatives to attributing a
gain to higher-order structure. Our MMR baseline shares the candidate set,
renderer, reader, and input limit with the learned selectors. A
size-constrained version additionally supports like-for-like retrieval
diagnostics.

\subsection{Learned interactions and optimization of learned scores}
Higher-order factorization machines provide low-rank ANOVA interactions
and efficient recurrences for training them~\cite{blondel2016hofm}.
Deep Sets represents permutation-invariant functions through pooled member
features~\cite{zaheer2017deepsets}. We adapt these standard forms to
query-conditioned evidence sets. Explicit fourth-order terms represent
interactions among distinct members; they do not by themselves identify
semantic relations or establish that a four-hop problem requires an
order-four model.

The mismatch between training and search is not a new observation in
machine learning. Beam-search optimization trains sequence scores on
search violations~\cite{wiseman2016bso}. Conservative objective models address
the tendency of optimization to exploit overestimated regions of a learned
objective~\cite{trabucco2021com}. Our intervention collects discrete
evidence sets encountered by the existing selector and trains preferences
from support annotations. It is neither a reproduction of sequence-level
BSO nor the conservative objective in COMs, and inherits no guarantee from
either. Their relevance is the need to evaluate scores where optimization
will use them.

\subsection{Granular and hypergraph representations}
Granular-ball computing groups observations at different scales
\cite{xia2019granularball}. In retrieval, Cross-Granularity Hypergraph RAG
uses entity--passage incidence and query-conditioned diffusion
\cite{wang2026crossgranularity}; HyperGraphRAG retrieves extracted n-ary
relations together with text~\cite{luo2025hypergraphrag}. These works precede
general claims about combining structure with retrieval. We use the full
name HyperGranular-RAG to avoid the shared HGRAG acronym, and do not equate
its old lexical facet hyperedges with factual relations or learned ANOVA
factors.

In the learned-selector study, a granular ball is an index node over
query-conditioned scoring vectors. The maximum-inner-product bound is a
standard center-plus-radius construction~\cite{ram2012tree}. K-means receives
the same vectors, bound, budget, and correctness requirements. Exact
agreement with flat search establishes preservation, not an independent
retrieval gain. Actual runtime determines whether the index is useful at
the tested candidate count.

\section{Task, evidence units, and information boundaries}
\subsection{Retrieval over the development paragraph corpus}
HotpotQA and MuSiQue provide multi-hop questions with supporting evidence
\cite{yang2018hotpotqa,trivedi2022musique}. We reuse 4,000 previously exposed
questions and reconstruct whole paragraphs from their original sentence
sources, retaining sentence-to-paragraph mappings. The resulting per-dataset
corpora contain 9,791 HotpotQA and 31,130 MuSiQue blocks. A fixed BGE-large-en-v1.5~\cite{xiao2023bge}
encoder supplies normalized 1,024-dimensional query and paragraph vectors,
and the top 128 paragraphs by dense score are shared by every selector.
Evaluation never inserts a paragraph because it has a support annotation.

This corpus is broader than each question's original closed candidate set,
but it is not full Wikipedia. There is no ANN approximation in this
candidate stage, and the corpora are assembled from the exposed benchmark
materials. Source availability and the corpus construction therefore limit
external validity. The BGE encoder uses its existing 512-token truncation
rule: 16 HotpotQA and 8 MuSiQue blocks exceed that limit. The reader receives
whole selected blocks under a separate input budget; encoder truncation and
reader feasibility are different operations.

Table~\ref{tab:roles} distinguishes training, selection, and QA roles.
The group split contains 3,166 FIT questions, 738 TUNE questions, and 96
historical-development questions. FIT has 3,165 groups. Question groups are
not split across roles, but the roles are not document-disjoint: FIT and TUNE
share 57 HotpotQA and 4,753 MuSiQue source documents. All are exposed
development material. In particular, excluding the answer-evaluation panel
from checkpoint selection does not retrospectively turn it into an
independent test set.
\begin{table}[t]\centering\small
\caption{Data roles (HotpotQA/MuSiQue question counts). MINE is within FIT;
DEV\_SELECT and QA are disjoint subsets of TUNE. All are exposed development
material. Repeated seeds reuse these same questions.}
\label{tab:roles}
\begin{tabular}{lrrl}\toprule
Role & HotpotQA & MuSiQue & Use\\\midrule
FIT & 748 & 2,418 & Supervised training\\
TUNE & 196 & 542 & Static model development\\
Historical DEV & 56 & 40 & Earlier interface development\\
MINE & 256 & 256 & Search-state collection\\
DEV\_SELECT & 64 & 64 & Continuation checkpoint choice\\
QA & 64 & 64 & Fixed reader evaluation\\\bottomrule
\end{tabular}\end{table}


\subsection{Feasibility and support labels}
For question $q$, let $V_q$ be the dense candidate list and $S\subseteq V_q$
a set of paragraph indices. Blocks are rendered with their titles and
original text, in BGE order. Define
\begin{equation}\label{eq:feasible}
\mathcal F_q=\{S\subseteq V_q:|S|\leq6,\;T(q,S)\leq1024\},
\end{equation}
where $T$ is the actual tokenizer length of the complete reader input,
including chat framing, question, and protocol instructions. We neither
truncate individual selected paragraphs nor assume that summing separately
tokenized lengths exactly predicts $T$. The cardinality and token limits
are both enforced. A long block may be infeasible even when a cardinality
slot remains available.

Let $G_q$ be the set of known support annotations, and let $C(S)$ be the
annotations recovered through the original source mapping. We define
\begin{equation}\label{eq:labels}
\operatorname{cov}(S)=\frac{|C(S)\cap G_q|}{|G_q|},\qquad
\operatorname{full}(S)=\mathbf1[G_q\subseteq C(S)].
\end{equation}
Full support is stricter than mean coverage. Neither is a complete label of
semantic usefulness: an unannotated paragraph can provide an alternative
derivation, and a reader can fail despite seeing all annotated support.
Targets supervise FIT sets and evaluate outputs; they are not inputs to the
deployment scorer, candidate ranking, or token feasibility check.

On all 738 TUNE questions, Top128 contains all known support for 190/196
HotpotQA and 367/542 MuSiQue questions. Known support fits the six-block and
input-token constraints for all 196 HotpotQA questions and 533/542 MuSiQue
questions, considered from their source blocks. These are separate
diagnostics: support may fit yet be absent from Top128. We retain the nine
budget-infeasible MuSiQue questions and every unreachable query in
unconditional results.

\section{Scoring sets and searching under the reader budget}
\subsection{Common representation and model families}
Let $e_q,e_i\in\mathbb R^{1024}$ be normalized embeddings and $l_i$ the
standalone paragraph token length. Each query--paragraph feature vector
concatenates $e_q,e_i,e_q\odot e_i,|e_q-e_i|$, dense cosine similarity,
and $\log(1+l_i)$. The last two features use FIT normalization statistics.
A linear layer, GELU, and layer normalization produce $h_i\in\mathbb R^{128}$;
a query network produces $h_q\in\mathbb R^{128}$. No answer, query ID, source
support identity, or rank label enters these features.

Unary outputs $a_{it}$ and query-conditioned cardinality biases $b_{qt}(k)$
are shared by the two prediction heads $t$: full support and coverage.
For interaction order $H$, factors $v_i\in\mathbb R^{32}$ use a linear
projection, layer normalization, and $\tanh$. Query weights
$w_{qth}\in\mathbb R^{32}$ score order $h$. The set logit is
\begin{equation}\label{eq:hofm}
g_{Ht}(S)=b_{qt}(|S|)+\frac{1}{6}\sum_{i\in S}a_{it}
+\sum_{h=2}^{H}\frac{w_{qth}^{\top}e_h(S)}{32\binom6h},
\end{equation}
where
\begin{equation}
e_h(S)=\sum_{i_1<\cdots<i_h,\;i_j\in S}v_{i_1}\odot\cdots\odot v_{i_h}.
\end{equation}
H1 has no interaction sum, H2 includes pairs, and H4 includes orders two
through four. The denominator is fixed by the maximum cardinality six,
not recomputed from $|S|$. An auxiliary coverage prediction is a separate
head; only the full-support logit ranks deployed sets.

The general-set control retains unary and cardinality outputs and adds
$\rho([\sum_{i\in S}h_i/6,h_q])$, a two-layer MLP with hidden width 192 and
two outputs. This is a DeepSets-style architecture, not a reproduction of
an LLM-based set selector. Parameter counts are comparable but not identical:
H2 has 670,640, H4 687,152, and DeepSets 707,922 parameters. H1 registers
662,384 parameters, of which 4,192 interaction-projection parameters do not
affect its score. Its effective scored parameter count is 658,192. We do
not conceal that inactive branch as useful first-order capacity.

\subsection{Static training and its evaluation target}
The initial training collection includes complete known-support sets,
sets with one known support removed, same-size replacements,
redundant additions, Dense-K6/MMR-K6, and bounded random examples.
FIT-only reference construction may append missing support blocks and
matched distractors to form supervised sets. That injection is limited to
static training examples. Actual evaluation and search mining always use
the original Top128 candidates.

The base objective combines binary cross-entropy for full support,
$0.5$ times coverage MSE after a sigmoid, $0.25$ times pointwise support
classification, and $0.25$ times a softplus ranking term for constructed
same-cardinality complete/incomplete pairs. Pointwise support labels again
mean annotation membership, not factual correctness. Initial checkpoints
were selected using original TUNE loss. Three initial training seeds
(1729, 2026, and 426) are retained. They share the same questions, inputs,
and answer-evaluation panel.

The constructed pair task is easier to control than deployment. It fixes
the alternatives presented to the model, often around known supports.
Search instead asks the model to choose among many feasible combinations.
Consequently, a high pair accuracy can be compatible with a low rate of
complete support in the selected output. We measure both endpoints rather
than treating the first as a substitute for the second.

\subsection{Actual Flat search and simple selectors}
The common learned search starts from the empty set. At each of six depths,
each surviving parent proposes its four highest-scoring feasible distinct
extensions, and the global beam retains four unique sets. All visited sets,
including those not retained for further expansion, remain eligible for
final selection. A feasible Dense-K6 set is also included. The final choice
maximizes the deployed full-support logit, with fewer actual input tokens
and then stable block identities breaking ties. Search does not terminate
merely because the next score increment is negative. This procedure is
exact with respect to each tested extension ranking, but the beam does not
enumerate all feasible sets and is not globally optimal.

Dense-budget greedily considers paragraphs in their initial dense order,
accepting a block when the complete input still fits. MMR-budget repeatedly
chooses the largest $0.7\,e_q^\top e_i-0.3\max_{j\in S}e_i^\top e_j$ among
remaining candidates, accepting it if feasible. Both skip overlong additions
and can retain more than six short blocks. Their K6 versions impose the
same cardinality as learned search. All selected contexts are finally
rendered in original BGE order, so MMR's selection sequence does not grant
it a different placement policy.

\subsection{What exact ball search can guarantee}
The ANOVA recurrence updates $e_h$ by descending order,
$e_h\leftarrow e_h+v_i\odot e_{h-1}$, with $e_0=\mathbf1$.
This is the standard higher-order factorization computation
\cite{blondel2016hofm}. For fixed parent $S$, the score of adding $i$ can be
written
\begin{equation}
g_H(S\cup\{i\})=o_q(S)+\theta_q(S)^\top z_i,
\end{equation}
where $z_i=[a_i/6,v_i]$ and the offset includes the parent score and
cardinality-bias change. Thus the conditional extension ranking is a maximum
inner-product problem in a learned query-specific space.

For a group with center $c_B$ and maximum Euclidean radius
$r_B=\max_{i\in B}\|z_i-c_B\|_2$, Cauchy--Schwarz gives
\begin{equation}\label{eq:bound}
\max_{i\in B}\theta^\top z_i\leq\theta^\top c_B+\|\theta\|_2r_B.
\end{equation}
The proof is immediate from $z_i=c_B+(z_i-c_B)$ and the definition of
$r_B$. An average cosine radius would not justify this bound. The
granular-ball index recursively partitions the scoring vectors with
farthest pivots; the k-means control uses the same space and bound with
iterated two-center assignments. Both preserve every original member and
perform the same token feasibility checks. They prune only groups unable
to improve the current feasible extension threshold, retaining ties and
numerical boundary cases.

Equation~\ref{eq:bound} is a standard search bound~\cite{ram2012tree}, not an
answer-quality theorem. Even exact preservation offers no rescue when the
score chooses an incomplete set. Nor does pruning prove a speedup: factor
construction, per-query indexing, traversal, and tokenizer work all contribute
to latency. These costs are measured separately in Section~\ref{sec:cost}.

\section{A bounded search-aligned correction}
\subsection{Mining the shared error distribution}
We select 512 FIT questions by deterministic group hash, with 256 per
dataset and proportional allocation to existing question types. Every
architecture searches these questions using its current score. The mined
pool merges final sets, high-scoring partial frontiers, Dense-K6/MMR-K6,
and at most eight deterministic feasible replacement attempts per model
trajectory. The merged, deduplicated pool is capped at 64 sets per question.
Model outputs are labeled from actual source identities, and are not
presumed incorrect merely because they were selected by a model.

All four architectures receive the same merged pool. No method receives
exclusive teacher calls, stronger candidates, or a different reader.
Unreachable queries still supply differences in partial annotated coverage;
we neither delete them nor invent a complete positive set. A possible second
round re-mines the current selected models rather than merely replaying the
first pool. The entire intervention is capped at two rounds.

The primary-seed pools retain 64 sets for each of the 512 queries in each
round. A missing preference is not fabricated: round one has no eligible
final-set pair for 171/256 HotpotQA and 115/256 MuSiQue queries, and no
frontier pair for 17/256 and 42/256. In round two these counts are 196/256,
135/256, 19/256, and 59/256, respectively. Pool capacity prioritizes model
finals and high-scoring frontiers, so attempted replacements need not survive
the cap. This limits the final-pair supervision actually available, even
though every query continues to contribute base observations.

\subsection{Preferences that train the deployed score}
Define the training preference $u(S)=0.5\operatorname{full}(S)+
0.5\operatorname{cov}(S)$. For $u(S^+)>u(S^-)$, use
\begin{equation}\label{eq:pair}
\ell=\log\!\left(1+\exp\{0.2[u(S^+)-u(S^-)]
-[g(S^+)-g(S^-)]\}\right).
\end{equation}
Final-set comparisons prefer equal cardinality and token-length ratios in
$[0.8,1.25]$, with same-cardinality fallback when no near-length pairs exist.
Frontier comparisons prefer the same parent and depth; same-depth comparisons
are a fallback when those pairs provide no preference. Missing pair
categories are masked and counted. The deployed logit $g$, rather than only
the auxiliary coverage head, receives gradients even when both partial
sets have zero full-support labels.

The total loss is
\begin{equation}
L=L_{\rm base}+0.5L_{\rm final}+0.5L_{\rm frontier}.
\end{equation}
Each query supplies four static and four mined base observations and up to
one sampled pair from each added category. Base and preference terms are
normalized within query and then across queries. More enumerated states
therefore do not automatically increase a query's training weight.
The preference is deliberately an annotation utility, not a reader reward.
No generated answer F1 is used to form pairs or choose a checkpoint.

\subsection{Controls and checkpoint selection}
Continuation starts from the corresponding original checkpoint with AdamW,
learning rate $3\times10^{-4}$, weight decay $10^{-4}$, query batch eight,
and gradient clipping at one. Each round permits eight epochs over the same
512-query mined FIT subset. H4 static replay uses eight static sets per
query, the original loss, and the same optimizer-update budget. Its reported
checkpoint matches the selected aligned H4 checkpoint's update count,
including the corresponding preceding-round starting point when applicable.
It is not independently chosen for favorable QA performance.
Equal updates isolate additional optimizer steps, not all training
computation: preference pairs add forward and backward evaluations.

A disjoint-from-QA checkpoint panel contains 128 TUNE questions, 64 per
dataset. At epochs 2, 4, and 8 we run actual selection. The fixed
lexicographic criterion maximizes dataset-equal full support, then coverage,
then minimizes original development loss, with earlier checkpoints breaking
remaining ties. The fixed natural-QA panel is excluded from this decision.
Because all materials have been exposed during development, the study is
still exploratory.

Both seeds (1729 and 2026) completed two rounds. The second seed was
triggered by selection-panel improvement before any new QA score was
observed; all four architectures received the same continuation.

The intervention jointly changes the example distribution, the added
ranking loss, and checkpoint selection (actual selected support rather than
the original TUNE loss). It therefore evaluates this training procedure as
a whole. Neither the replay control nor the two-round design identifies
the ranking loss's isolated causal effect.

\section{Evaluation and reproducibility}
\subsection{Reader and canonical scores}
The reader is the original Qwen2.5-3B-Instruct checkpoint~\cite{qwen2024release} in FP16, without
the relation-front-end adapter. Its P2 interface requests a single answer
field, uses constrained decoding and greedy generation, and permits 128
new tokens. Every method shares the tokenizer, chat framing, original-text
renderer, precision, and input limit. A public protocol improvement is not
attributed to granular balls or higher-order scoring.

The main panel has 128 questions: 48 HotpotQA bridge and 16 comparison
questions, and 32 MuSiQue two-hop and 32 three/four-hop questions. Canonical
dataset-specific EM and token F1 score the parsed answer field, including
each dataset's normalization and alias conventions. We do not apply
HotpotQA-specific yes/no handling indiscriminately to MuSiQue, and do not
clean answers using the target. Invalid output remains a failure in the
full denominator. UNKNOWN, EOS, output caps, and format failures are
recorded separately.

The primary continuation comparison retains Dense, MMR, the original H4,
static H4 replay, and H1/H2/H4/DeepSets aligned models. A logical result
can reuse a prior generation only when full input token IDs, model,
precision, generation configuration, and constrained-decoding identity
match. Merely sharing evidence text does not satisfy that condition.
Cache reuse is reported as reuse, not independent inference or a deployment
latency advantage.

\subsection{Interpretation of repeated comparisons}
Dataset means are reported separately and combined with equal dataset
weight. The same questions appear across methods and training seeds.
Seed standard deviations describe training variation, not confidence
intervals over independent questions. The continuation does not introduce
new confirmation p-values, and no nonsignificant or close comparison is
called equivalent. We retain failures, unreachable support, and contexts
that cannot accommodate all reference evidence. Conditional diagnostics
are labeled and do not replace unconditional means.

Historical sentence-completion results in Section~\ref{sec:history} use
their existing canonical rescoring and original paired-bootstrap intervals.
That additive re-evaluation fixes predictions rather than generating new
ones and is not a new independent confirmation. Those intervals are not
transferred to the current learned-selector comparisons.

\section{Results: from proxy fit to actual selection}
\subsection{The original mismatch}
On constructed same-size complete/incomplete pairs, H4 achieves accuracy
0.9645 on HotpotQA and 0.9464 on MuSiQue. Yet its actual selected support
completeness on all 738 TUNE queries is 0.7092 and 0.2232, below Dense-K6's
0.8214 and 0.2565. Table~\ref{tab:static} includes the lower-order and general
set alternatives. H4's favorable pair metric therefore cannot establish a
practical selection advantage.

\begin{table}[t]\centering
\caption{Original static models (seed 1729) on the 738-query TUNE boundary. Pair accuracy
uses constructed same-cardinality complete/incomplete pairs. Full support
uses the actual chosen set. H/M denote HotpotQA/MuSiQue.}
\label{tab:static}
\begin{tabular}{lrrrr}\toprule
Method & Pair H & Pair M & Full H & Full M\\\midrule
Dense-K6 & -- & -- & .8214 & .2565\\
H1 & .9475 & .8978 & .6531 & .2454\\
H2 & .9638 & .9422 & .6429 & .2325\\
DeepSets & .9620 & .9482 & .6990 & .2306\\
H4 & .9645 & .9464 & .7092 & .2232\\\bottomrule
\end{tabular}
\end{table}

The loss is visible at the query level. Relative to Dense-K6, H4 loses full
support on 27 HotpotQA and 42 MuSiQue queries while recovering it on only
5 and 24. The selected set scores are at least as high as the included
Dense-K6 scores, so this is not caused by omitting that fallback from the
score comparison. Mean sigmoid outputs on selected sets are 0.9920 and
0.9667, far above their observed full-support frequencies. We describe these
as model outputs, not calibrated probabilities. The scorer is optimizing
its learned criterion successfully while that criterion can favor poorer
annotated evidence.

Additional controls limit the interpretation of high pair accuracy.
On exact-length and same-cardinality matched comparisons, H4 accuracy is
0.8333/0.7878, compared with DeepSets 0.9524/0.8525. These are smaller
diagnostic subsets: 65 HotpotQA queries with 84 pairs and 185 MuSiQue queries
with 278 pairs. They do not replace the full selection boundary, but weaken
an explanation based on uniquely superior higher-order discrimination.
High-order parameters receive nonzero gradients and contribute nonzero
terms, so their existence in the model is not merely nominal. Their
activation still does not establish their necessity.

\subsection{Search-aligned development results}
\begin{table}[t]\centering\small
\caption{Checkpoint-panel support completeness before and after search alignment (64 questions per dataset). Selection uses this panel, not QA.}
\label{tab:aligned}
\begin{tabular}{llrrrr}\toprule
Seed & Model & Static H & Aligned H & Static M & Aligned M\\\midrule
1729 & H1 & 0.5938 & 0.7500 & 0.2656 & 0.2656\\
1729 & H2 & 0.6406 & 0.7500 & 0.2344 & 0.2500\\
1729 & DeepSets & 0.6719 & 0.7500 & 0.2188 & 0.2656\\
1729 & H4 & 0.6719 & 0.7500 & 0.2344 & 0.2812\\
2026 & H1 & 0.6094 & 0.7656 & 0.2344 & 0.2656\\
2026 & H2 & 0.6406 & 0.7188 & 0.2500 & 0.2500\\
2026 & DeepSets & 0.6250 & 0.7812 & 0.2344 & 0.2656\\
2026 & H4 & 0.6719 & 0.7188 & 0.2344 & 0.2500\\
\bottomrule\end{tabular}\end{table}

Table~\ref{tab:aligned} reports DEV\_SELECT, the panel used to choose these
checkpoints. For seed 1729, H4's dataset-equal annotated support completeness
rises from 0.4531 to 0.5156. H2 and DeepSets show the same increase, 0.0625;
H1 increases by 0.0781. These checkpoint-selected effects cannot establish
external generalization or a fourth-order advantage. All eight reported
aligned checkpoints come from round one; the completed second rounds do
not improve the best selection criterion. QA outcomes on a separate subset
of the same exposed TUNE pool follow in Section~\ref{sec:answers}.


\section{Results: natural answers and computational cost}\label{sec:answers}
\subsection{The fixed natural-QA panel}
\begin{table}[t]\centering\small
\caption{Same QA panel (64 questions per dataset), seed 1729. Full is annotated support completeness, with counts out of 128. F1 is dataset-equal canonical answer F1; $\Delta$F1 is relative to MMR. Dense/MMR share the token ceiling but may use more than six blocks.}
\label{tab:qa}
\begin{tabular}{lrrr}\toprule
Method & Full (count) & F1 & $\Delta$F1\\\midrule
Dense & 0.5000 (64) & 0.4055 & -0.0248\\
MMR & 0.4766 (61) & 0.4303 & +0.0000\\
Static H4 & 0.4297 (55) & 0.4181 & -0.0122\\
H4 replay & 0.4297 (55) & 0.4046 & -0.0257\\
H1 aligned & 0.4766 (61) & 0.4065 & -0.0238\\
H2 aligned & 0.4375 (56) & 0.4057 & -0.0245\\
DeepSets aligned & 0.4688 (60) & 0.4090 & -0.0213\\
H4 aligned & 0.4453 (57) & 0.3766 & -0.0537\\
\bottomrule\end{tabular}\end{table}

Table~\ref{tab:qa} puts annotated support and answers on the same QA panel.
For seed 1729, H4 completeness changes from 55/128 (0.4297) to 57/128
(0.4453), while F1 falls from 0.4181 to 0.3766. The paired support counts
are four gains and two losses; answer F1 improves on 10 questions, falls on
14, and is unchanged on 104. These are marginal changes, not evidence that
adding support causes an answer loss. For seed 2026, H4 completeness changes
from 53/128 to 57/128 and F1 from 0.4304 to 0.3831. Aligned H1, H2, and
DeepSets score 0.4126, 0.4173, and 0.3521 in that seed. Every method retains
all 128 exposed questions. Neither seed establishes an F1 advantage over MMR.

Relative to each corresponding static checkpoint, the primary-seed mean F1 changes are H1 $+0.0326$, H2 $+0.0060$, DeepSets $-0.0033$, H4 $-0.0415$.  For seed 2026, the corresponding changes are H1 $-0.0151$, H2 $+0.0051$, DeepSets $-0.0505$, H4 $-0.0473$. The static comparisons reuse the same original questions and predictions; no additional generation or independent replication is implied.

\begin{table}[t]\centering\small
\caption{Primary-seed QA-panel selection and output diagnostics, equal dataset weight over all 128 questions. Full is annotated support completeness; EM is canonical answer exact match. Input includes instructions.}
\label{tab:qa-details}
\begin{tabular}{lrrrrr}\toprule
Method & Full & Coverage & Blocks & Input tokens & EM\\\midrule
Dense & 0.5000 & 0.7096 & 6.95 & 1009.9 & 0.3359\\
MMR & 0.4766 & 0.7025 & 7.25 & 1007.3 & 0.3672\\
Static H4 & 0.4297 & 0.6608 & 5.91 & 916.9 & 0.3594\\
H4 replay & 0.4297 & 0.6647 & 5.91 & 900.9 & 0.3359\\
H1 aligned & 0.4766 & 0.6960 & 4.45 & 696.2 & 0.3359\\
H2 aligned & 0.4375 & 0.6667 & 5.53 & 789.7 & 0.3516\\
DeepSets aligned & 0.4688 & 0.6914 & 4.43 & 693.4 & 0.3516\\
H4 aligned & 0.4453 & 0.6803 & 5.70 & 842.2 & 0.2969\\
\bottomrule\end{tabular}\end{table}


Table~\ref{tab:qa-details} separates the cardinality actually used, the full
tokenized input, annotated support, and exact answer match. The nominal
1,024-token cap is shared, not a claim of identical input lengths. Dense and
MMR may use more short paragraphs than the learned six-block selectors.
Consequently, the answer comparison measures complete selection pipelines
under a common token ceiling, not an isolated model-order effect at identical
cardinality or text length. Dense-K6 supplies the separate matched-cardinality
support comparison in Table~\ref{tab:static}.

The original three-seed results provide context without being pooled with
the continuation as new replications. Dataset-equal mean F1 is 0.4048 for
H1, 0.4109 for H2, 0.4041 for DeepSets, and 0.4171 for H4. Dense and MMR
score 0.4055 and 0.4303. H4's seed range is 0.4027--0.4304, and H1/H2 exceed
it for seed 426. Thus a favorable single checkpoint does not establish a
general interaction-order advantage. The fixed reader's response can
change when an annotation-equivalent set contains different context, so
support and answer comparisons must both remain visible.

Earlier 64-query token-budget sensitivity also remains a separate boundary.
At 512/1,024/2,048 tokens, H4's equal-dataset F1 is
0.4020/0.4649/0.4857, while MMR's is 0.4281/0.4489/0.3737.
These measurements use the same smaller subset across token budgets, not
the full 128-question panel. They suggest dependence on context and
selection policy but do not license choosing 2,048 tokens post hoc for the
present main result. The continuation retains the prescribed 1,024-token
budget and does not run a new budget grid.

\subsection{Search correctness and latency}\label{sec:cost}
Exact granular-ball and k-means search agree with Flat on 1,476 comparisons
over all 738 TUNE queries and 30,956 tested parent states. A bounded
implementation improvement that skips demonstrably inferior exact leaf
scores was subsequently tested on a fixed 32+32-query subset, producing
128 comparisons and 2,688 parent states with no mismatch. Table~\ref{tab:cost}
reports that latter timing boundary. The initial full-boundary timings are
also retained in the reproducibility materials rather than overwritten.

\begin{table}[t]\centering
\caption{Existing fixed-subset mean selection time in seconds per query,
including network, per-query index construction, and tokenizer feasibility.
These component timings precede the search-aligned training intervention.}
\label{tab:cost}
\begin{tabular}{lrrr}\toprule
Dataset (queries) & Flat & Granular ball & K-means\\\midrule
HotpotQA (32) & .33459 & .42557 & .42268\\
MuSiQue (32) & .54612 & .65282 & .64392\\\bottomrule
\end{tabular}
\end{table}

Granular balls reduce inner-product evaluations by approximately 66\% on
HotpotQA and 44\% on MuSiQue, but their measured latency remains 27.2\%
and 19.5\% above Flat. Both tree types incur construction and traversal
overhead and can require more token checks before establishing feasible
incumbents. All candidate embeddings and small-network factors are still
computed. Some CPU measurements overlap GPU reader work, which precludes
claims of isolated device-independent speed. The evidence supports exact
preservation at this boundary, not acceleration.

Replacing each leaf by one permanent representative is a different,
approximate method. It changes all 738 selected sets and reduces complete
support to 0.3469/0.0535. This diagnostic explains why preserving every raw
member matters here. It does not refute every approximate grouping method,
and no additional reader outputs were generated for that control.

\subsection{Model calls and experiment cost}
The continuation records 4,513.07 GPU process-wall seconds, 6,997.06 measured research CPU seconds, and 986 real reader calls, including eight uncached checks. It produces 1,536 logical QA rows with 558 cache hits. New generations total 761,847 input tokens and 9,460 output tokens. No comparison row has an invalid answer field or reaches the output limit. Typesetting CPU is accounted separately in the accompanying ledger. Older unmeasured CPU use remains unknown; phase accounting is not reconstructed whole-system billing. No paid compute is used.

The prior static-selector study used 4,350.28 GPU process-wall seconds,
10,541.75 measured CPU seconds, and 2,287 real reader calls including eight
uncached checks. Its 3,520 logical rows contain 1,241 cache hits. Those rows
are neither 3,520 independent examples nor 3,520 fresh generations.

\section{Context from earlier evidence-completion studies}\label{sec:history}
The earlier static-q25 system used sentence units, compact MiniLM geometry,
query-term facet relations, and at most four inserted units after ten
protected dense entries. It retained a final Top-20 entry budget under a
separate 4,096-token input cap and used Qwen2.5-1.5B-Instruct. Its facet
proposal was scored once relative to seed groups; it was not the dynamically
recomputed learned set score of this article. Table~\ref{tab:history}
preserves canonical effects and recorded intervals under their original
boundaries.

\begin{table}[t]\centering\small
\caption{Historical sentence-completion evidence, kept separate from the
paragraph-selector experiments. Effects are answer-F1 differences, with
recorded 95\% paired-bootstrap intervals. Joint studies use equal dataset
weight. Unmarked component/BGE/placement rows contain 2,500 queries on their
respective boundaries; Gemma transfer contains 4,000. Each contrast retains
its original full denominator; different boundaries are not pooled.}
\label{tab:history}
\begin{tabular}{p{.52\linewidth}rr}\toprule
Comparison & $\Delta$F1 & 95\% interval\\\midrule
Compact minus Dense, HotpotQA (1,000) & .01478 & [.00020,.02988]\\
Compact minus Dense, MuSiQue (3,000) & .01140 & [.00450,.01835]\\
Compact minus Dense, joint (2,500) & .01357 & [.00491,.02233]\\
Compact Full minus NoFacet & .01336 & [.00341,.02343]\\
Original compact Full minus BGE & -.03998 & [-.05393,-.02621]\\
BGE sidecar Protected minus BGE & -.00256 & [-.00998,.00458]\\
BGE-native Protected minus BGE & -.00305 & [-.00688,.00063]\\
Gemma transfer, completion minus Dense & .00516 & [-.00262,.01295]\\
Cross-space Protected minus Unprotected & .01122 & [.00129,.02104]\\\bottomrule
\end{tabular}
\end{table}

Three compact-backbone answer improvements are supported within those
closed-candidate conditions. The specified facet selector also exceeds its
centroid-only control, but that comparison does not isolate every possible
coverage or diversity heuristic. The original compact Full system is worse
than BGE. Neither its cross-space sidecar nor a BGE-native reconstruction
establishes an increment over BGE, and the Gemma transfer interval crosses
zero. These unresolved effects are not evidence of equivalence or of
universal ineffectiveness.

The original Full versus NoProtection ablation remains unresolved
($\Delta$F1 $=0.00354$, recorded interval $[-0.00675,0.01389]$). It does not
establish the independent necessity of the old protection component. This
ablation and the later fixed-content ordering comparison test different
interventions and must not be substituted for one another.

The cross-space placement comparison holds prompt-visible evidence content
fixed and changes its ordering. With outside-Top20 insertions $I$, $|I|=m\leq4$,
both policies contain $I\cup\{d_1,\ldots,d_{20-m}\}$, and the completed
position audit verified the visible text and truncation conditions. Its
positive effect therefore describes the tested ordering/placement policy
on all original queries, not an internal attention mechanism or a favorable
post hoc subset. A policy can improve how a reader uses a fixed set without
establishing that the set improves BGE's own evidence.

These results motivate studying evidence selection more directly, but they
cannot validate the new learned models. Conversely, a successful search
correction would not retroactively prove that the old granular-ball or
facet construction caused the compact gains. The two studies differ in
candidate universe, unit granularity, generator, prompt, and budgets.

\section{Discussion}
\subsection{A score must be evaluated under its optimizer}
The constructed-set results show that the models learn a useful supervised
distinction, but the deployment results identify a different requirement.
Optimization changes the distribution of evaluated sets. A model may
recognize a supplied missing-support contrast while assigning excessive
scores to other combinations that its search discovers. Including Dense-K6
as a feasible candidate makes the failure particularly informative: the
search can reject an available, more complete set because the deployed
criterion ranks it lower.

The continuation targets this discrepancy directly. Its partial-state
preferences can change the actual search order even when no set in a
comparison is fully complete. Static replay matches additional optimizer
steps, while the intervention also changes checkpoint selection and training
computation. Low-order and general-set controls
test whether explicit fourth-order interactions are responsible for any
gain. An improvement shared by those controls is evidence for training
alignment within this setup, not evidence for a unique granular or
hypergraph mechanism.

\subsection{Evidence quality and reader quality remain distinct}
Support annotation offers inexpensive supervision with interpretable labels,
but it is an incomplete account of useful context. Alternative evidence,
entity disambiguation, distracting text, and the reader's own knowledge can
all change an answer without changing known-support completeness. Conversely,
complete annotated support does not guarantee that the reader will combine
the facts correctly. Reader protocol and constrained-output changes belong
to the public interface and must not be attributed to the evidence model.

This distinction also limits a negative interpretation. A model's failure
to improve F1 on one exposed panel does not prove that its evidence is
useless on all readers. Equally, a positive point estimate cannot establish
generalization after repeated development. The appropriate conclusion is
the effect actually measured under the shared candidate, budget, and reader
conditions, together with the controls that explain or fail to explain it.

\subsection{Indexing is a systems hypothesis}
Exact indexing asks whether the same answer to a search problem can be
obtained more cheaply. It does not add evidence-selection capacity when it
preserves the scorer and beam. At only 128 candidates, vectorized Flat
search has small arithmetic cost, while index traversal and actual-token
feasibility remain visible. The ball and k-means results therefore favor
Flat in this implementation. A larger corpus or another feasibility regime
might change the balance, but no such result is claimed here.

The practical lesson is to evaluate a sequence of interfaces rather than
one attractive metric: supervision must train the deployed score; the
score must choose useful evidence; the reader must use it; and any
acceleration must include the costs it introduces. This decomposition is
applicable as an evaluation strategy beyond the present model family,
although its empirical outcomes are bounded to the tested system.

\section{Limitations}
All learned-selector and continuation results use exposed development
materials. There is no new independent confirmation, full-Wikipedia run,
new language, or external-method deployment. Document overlap between FIT
and TUNE prevents unseen-document claims. The answer panel is small and
balanced by design rather than representative of an arbitrary deployment
distribution. Training-seed variation cannot replace query-level external
validation, and resource-bounded continuation does not rule out all future
learning objectives.

Known supports are incomplete semantic labels. Some required blocks are
unreachable or too long, and the encoder itself truncates a small number of
paragraphs. We report these boundaries rather than labeling every lost or
unannotated passage as noise. The use of a single 3B reader and a fixed
structured-output protocol limits generator conclusions. Earlier generator
transfer remains unresolved under its own conditions.

The study does not reproduce SetR, RECOMP, Beam Retrieval, or a large
graph-RAG system under matched conditions. It supplies simple and internal
learned controls, not a comprehensive leaderboard. The historical facet
comparison leaves ordinary clustering and direct selection attribution
unresolved. Standard ANOVA algebra and maximum-radius bounds are included
to define the evaluated algorithms; they are not new theoretical
contributions or guarantees about answer F1.

\section{Conclusion}
Budgeted evidence selection should be judged where its score is used.
In this study, strong constructed-set prediction coexists with losses of
annotated support during search, and higher-order modeling does not by
itself establish an answer-quality advantage over MMR. A bounded continuation
improves checkpoint-panel support selection without changing the reader or
candidate universe. On the separate, fixed QA panel, primary-seed H4 gains
two net complete-support questions (55 to 57 of 128), yet answer F1 declines
from 0.4181 to 0.3766. H4 answer F1 declines in both training seeds and
every aligned architecture remains below MMR in F1 point estimate.
H2 improves slightly over its own static checkpoint in both seeds, while
H1 improves in only the primary seed. The answer effects therefore depend
on architecture and seed. Exact granular-ball and k-means indexes
preserve selection but provide no measured acceleration at 128 candidates.
Together with the separately reported historical completion and placement
results, the evidence supports an empirical account of the gap between
proxy objectives, selected context, answers, and cost. Stronger deployment
claims require new prospective evidence, not reinterpretation of the
development comparisons reported here.

\backmatter
\section*{Statements and declarations}
\paragraph{Data and code availability.}
Online Resource 1 provides anonymized per-query numeric results, a portable
script that reconstructs the QA tables and paired descriptive counts, and
data, model, and preprocessing specifications. The public benchmark and
model sources are identified there. Text, answers, model weights, and
inference caches are not redistributed. Table reconstruction runs without
these assets; a complete model rerun additionally requires source-data
access, model licenses, and the experiment implementation. The supplement
is a numerical reproducibility package, not a claim of a self-contained
one-command training and generation release.
\paragraph{AI assistance.}
OpenAI Codex contributed implementation, analysis, targeted literature
checking, and manuscript drafting and revision. This extends beyond
copy editing and includes generative editorial work.
\paragraph{Research ethics.}
This computational study uses the cited benchmark materials and existing
models. It did not recruit human participants or collect new personal data.
The supplementary release excludes original text and answer annotations.

\bibliography{references,additions}
\end{document}
